静电场既可以用电场强度描述，也可以用能量描述。对于电荷量为 $q$ 的试探电荷，电场力
$\vect F=q\vect E$ 是保守力，因此可以定义电势能 $U$。把电势能中与试探电荷成正比的部分
分离出来，定义电势
\begin{equation}
  \boxed{\varphi=\frac{U}{q}},
  \qquad U=q\varphi.
\end{equation}
电势是电场自身的标量属性；电势能则属于“电荷与电场组成的系统”，会随试探电荷的大小和
正负改变。

\subsection*{电势与电场}

保守力等于势能梯度的负值。由 $\vect F=-\nabla U$、$U=q\varphi$ 和
$\vect F=q\vect E$ 可得
\begin{equation}
  \boxed{\vect E=-\nabla\varphi}.
\end{equation}
在直角坐标中即
\begin{equation}
  E_x=-\frac{\partial\varphi}{\partial x},
  \qquad
  E_y=-\frac{\partial\varphi}{\partial y},
  \qquad
  E_z=-\frac{\partial\varphi}{\partial z}.
\end{equation}
反过来，试探电荷从 $A$ 运动到 $B$ 时，电场力做功为
\begin{equation}
  W_{A\to B}=q\int_A^B\vect E\cdot\dd\vect l
  =q(\varphi_A-\varphi_B),
\end{equation}
所以
\begin{equation}
  \boxed{\varphi_B-\varphi_A
  =-\int_A^B\vect E\cdot\dd\vect l}.
\end{equation}
静电场中该积分与路径无关。电场始终指向电势降低最快的方向。

\potentialGradientDiagram

电势零点可以任意选择。对空间中局域的有限电荷分布，通常取无穷远处电势为零；但无限长
直线或无限大平面的电势在无穷远处发散，只能选一个有限参考点，计算电势差而不能机械地令
$\varphi(\infty)=0$。

\subsection*{常见电荷分布的电势}

取无穷远为零电势，点电荷 $Q$ 在距它 $r$ 处产生的电势为
\begin{equation}
  \boxed{\varphi=\frac{1}{4\pi\varepsilon_0}\frac{Q}{r}}.
\end{equation}
电势是标量，满足代数叠加：
\begin{equation}
  \varphi=\frac{1}{4\pi\varepsilon_0}\sum_i\frac{Q_i}{r_i},
  \qquad
  \varphi=\frac{1}{4\pi\varepsilon_0}\int\frac{\dd q}{r}.
\end{equation}
复杂电荷分布常可先叠加电势，再利用 $\vect E=-\nabla\varphi$ 求电场，从而避免直接进行
矢量积分。

电偶极子的电偶极矩定义为 $\vect p=q\vect l$，方向由负电荷指向正电荷。若场点到偶极子
中心的距离 $r\gg l$，则
\begin{equation}
  \boxed{\varphi(\vect r)=
  \frac{1}{4\pi\varepsilon_0}\frac{\vect p\cdot\vect r}{r^3}
  =\frac{1}{4\pi\varepsilon_0}\frac{p\cos\theta}{r^2}}.
\end{equation}

半径为 $R$、总电荷为 $Q$ 的均匀带电薄球壳，其电势为
\begin{equation}
  \boxed{\varphi(r)=
  \begin{cases}
    \dfrac{Q}{4\pi\varepsilon_0r}, & r\ge R,\\[6pt]
    \dfrac{Q}{4\pi\varepsilon_0R}, & r<R.
  \end{cases}}
\end{equation}
球壳外等效于球心处的点电荷；球壳内电势为常量，因此电场为零。

\subsection*{等势面}

空间中电势相等的点组成等势面。它具有三个直接性质：
\begin{enumerate}[label=\arabic*.,leftmargin=2em]
  \item 电场线处处垂直于等势面；
  \item 不同电势的等势面不能相交；
  \item 当相邻等势面所代表的电势差相同时，等势面越密集，电场越强。
\end{enumerate}
电荷沿等势面移动时，电场力不做功。对于给定电势函数，等势面方程就是
$\varphi(x,y,z)=\varphi_0$。

\dipoleEquipotentialDiagram

\subsection*{电子伏特与微观能标}

一个带一个元电荷的粒子通过 $\SI{1}{V}$ 电势差所改变的能量定义为
\begin{equation}
  \boxed{\SI{1}{eV}=1.602176634\times10^{-19}\ \si{J}}.
\end{equation}
电子伏特是能量单位，不是电势单位。原子电离能通常为数电子伏特到数十电子伏特；核反应
常以 \si{MeV} 为单位；高能粒子物理常使用 \si{GeV}、\si{TeV}。热运动的典型能量约为
$k_BT$，令 $k_BT=\SI{1}{eV}$，可得
\begin{equation}
  T\approx1.16\times10^4\ \si{K}.
\end{equation}

\ionizationEnergyDiagram

带电粒子每经过一段加速电势差，能量便增加 $q\Delta\varphi$。大型粒子加速器用多级电场
不断提高粒子能量，因此常用 \si{GeV}、\si{TeV} 描述其能标。

\historyfigure[0.72\linewidth]{particle-accelerator-redraw.png}{大型粒子对撞机的环形隧道与加速装置}

图中的环形真空管道用于让粒子束反复循环；真正增加粒子能量的是沿轨道布置的加速电场，
磁场主要负责弯转和聚焦。粒子每通过一次有效加速电势差 $\Delta\varphi$，能量增加
$q\Delta\varphi$，经过许多次加速后累积到 \si{GeV} 或 \si{TeV} 能标。这说明电子伏特的定义
可以直接连接电势差与加速器中的粒子能量。

气体放电可作数量级估算：若原子典型电离能约为 \SI{10}{eV}，电子平均自由程为 $l$，则
电场应大致满足 $eEl\gtrsim\SI{10}{eV}$。常温常压气体中 $l$ 为亚微米到微米量级，所以
击穿场强通常处于 $10^6$ 至 $10^7\ \si{V.m^{-1}}$ 的数量级，实际值还取决于气体种类、
压强和电极形状。

\begin{center}
\begin{minipage}{0.92\linewidth}
  \begin{minipage}[c]{0.27\linewidth}
    \centering
    \includegraphics[width=\linewidth]{lightning-discharge-redraw.png}\par\smallskip
    {\small\color{TextGray}闪电形成的气体放电通道}
  \end{minipage}\hfill
  \begin{minipage}[c]{0.68\linewidth}
    图中明亮而分叉的部分是空气击穿后形成的导电通道，并不是电场线。击穿之前，自由电子在
    平均自由程 $l$ 内由电场获得的能量约为 $eEl$；当它足以引起碰撞电离时，电子数迅速倍增，
    放电通道才会发光并向前发展。因此图中的闪电通道与数量级条件
    $eEl\gtrsim\SI{10}{eV}$ 相对应。
  \end{minipage}
\end{minipage}
\end{center}

\subsection*{电荷系统的静电势能}

由 $N$ 个点电荷组成的系统，其总静电势能等于所有不同电荷对相互作用能之和：
\begin{equation}
  \boxed{U=\frac{1}{4\pi\varepsilon_0}
  \sum_{i<j}\frac{Q_iQ_j}{r_{ij}}}.
\end{equation}
也可写成
\begin{equation}
  \boxed{U=\frac12\sum_iQ_i\varphi_i^{(\mathrm{others})}},
\end{equation}
其中 $\varphi_i^{(\mathrm{others})}$ 不含 $Q_i$ 自身产生的电势。因子 $1/2$ 用来消除每一对
电荷被重复计算的问题。

对连续电荷分布，在自能有限且积分有意义时，推广为
\begin{equation}
  U=\frac12\int\rho\varphi\,\dd V.
\end{equation}
真空静电场的能量密度为
\begin{equation}
  \boxed{u_E=\frac12\varepsilon_0E^2},
\end{equation}
因而整个电场的能量为
\begin{equation}
  \boxed{U=\frac{\varepsilon_0}{2}
  \int_{\mathbb R^3}E^2\,\dd V}.
\end{equation}
对于点电荷，靠近点电荷处的积分发散；这反映了经典点电荷模型具有无穷自能。对于有限尺寸
的带电体，上式通常给出有限结果。

\subsection*{从做功定义电势差}

电势的绝对值依赖零点选择，电势差却不依赖。若把单位正电荷从 $A$ 缓慢移到 $B$，外力做功
等于电势能增加量，因此
\begin{equation}
 \varphi_B-\varphi_A=\frac{W_{\rm ext}}q=-\frac{W_{\rm field}}q.
\end{equation}
“缓慢”意味着过程中的动能变化可以忽略。若电荷自由运动，则应把电场力做功写进动能定理：
\begin{equation}
 \frac12mv_B^2-\frac12mv_A^2=q(\varphi_A-\varphi_B).
\end{equation}
正电荷由高电势向低电势加速，负电荷则相反；无论电荷正负，电势能降低量都转化为动能增加量。

静电场保守性的局部形式为 $\nabla\times\vect E=0$。在单连通区域中，它保证线积分与路径无关，
从而可以定义单值电势。反过来，由 $\vect E=-\nabla\varphi$ 自动得到零旋度。若存在随时间变化的
磁场，电场环流一般不再为零，单一静电势便不能完整描述电场；本讲讨论的电势关系限定在静电场。

沿一条给定方向 $s$，电场分量为
\begin{equation}
 E_s=-\frac{\dd\varphi}{\dd s}.
\end{equation}
因此电势--位置图的斜率直接给出电场。电势有极值的点满足 $\vect E=0$，但零场点未必稳定；
还需考察电势的二阶变化。对正试探电荷，势能 $q\varphi$ 的局部极小对应稳定平衡，局部极大
对应不稳定平衡；负电荷的判断相反。

\subsection*{由电荷求电势，再由电势求电场}

电势叠加是标量代数和，通常比直接叠加电场简单。对连续分布可写成
\begin{equation}
 \varphi(\vect r)=\frac1{4\pi\varepsilon_0}
 \int\frac{\rho(\vect r')}{|\vect r-\vect r'|}\,\dd V'.
\end{equation}
若电荷分布在有限区域内且观察距离远大于分布尺度，电势可按距离反幂展开：
\begin{equation}
 \varphi(\vect r)=\frac1{4\pi\varepsilon_0}
 \left(\frac Qr+\frac{\vect p\cdot\uvec r}{r^2}+\cdots\right).
\end{equation}
总电荷 $Q$ 非零时，最远处由单极项 $1/r$ 主导；若 $Q=0$ 而偶极矩 $\vect p$ 非零，则由
$1/r^2$ 的偶极项主导；更高阶四极项衰减更快。这解释了为什么中性分子远场仍可能产生电作用，
同时也说明高阶多极结构只在较近距离才明显。

电势满足 Poisson 方程
\begin{equation}
 \boxed{\nabla^2\varphi=-\frac\rho{\varepsilon_0}},
\end{equation}
无电荷区域满足 Laplace 方程 $\nabla^2\varphi=0$。一维平面对称问题中它退化为
$\dd^2\varphi/\dd x^2=-\rho/\varepsilon_0$；分段电荷分布可在各区分别积分，再用电势连续和
法向电场跃变条件连接。

\subsection*{等势面的几何信息}

等势面与电场正交可以由微分直接看出。沿等势面任意切向位移 $\dd\vect l_t$，有
\begin{equation}
 \dd\varphi=\nabla\varphi\cdot\dd\vect l_t
 =-\vect E\cdot\dd\vect l_t=0,
\end{equation}
故电场没有切向分量。若在一幅等势图上相邻曲线的电势差都为 $\Delta\varphi$，局部法向间距为
$\Delta n$，则场强近似为 $E\approx|\Delta\varphi|/\Delta n$。等势面密集处场强大，但只有
相邻面所代表的电势差相同，这种疏密比较才有意义。

无限长线电荷是电势零点选择的典型例子。由高斯定律
$E(r)=\lambda/(2\pi\varepsilon_0r)$，两点间电势差为
\begin{equation}
 \varphi(r_2)-\varphi(r_1)
 =-\frac\lambda{2\pi\varepsilon_0}\ln\frac{r_2}{r_1}.
\end{equation}
当 $r_2\to\infty$ 时对数发散，所以不能把无穷远取为零；选择任意有限参考半径并不会影响电场
或可观测的电势差。

\subsection*{电荷系统能量的建立过程}

把一组电荷从无穷远逐个搬到指定位置，可以得到系统势能。第一个电荷搬入时没有其他电荷，
不需克服相互作用；第二个电荷只与第一个相互作用；继续进行，每一对电荷恰好被计算一次，于是
得到 $\sum_{i<j}kQ_iQ_j/r_{ij}$。若改写为 $\frac12\sum_iQ_i\varphi_i^{(\rm others)}$，每一对
会在 $i$ 与 $j$ 两项中出现两次，所以必须乘 $1/2$。

连续分布的能量也可从场的观点理解。利用 $\rho=\varepsilon_0\nabla\cdot\vect E$、
$\vect E=-\nabla\varphi$ 和分部积分，在边界项消失时有
\begin{equation}
 \frac12\int\rho\varphi\,\dd V
 =\frac{\varepsilon_0}{2}\int E^2\,\dd V.
\end{equation}
左式把能量看作电荷之间的相互作用，右式把同一能量看作储存在整个空间的电场中。场能密度
处处非负，但选定零点后的两个异号点电荷相互作用能可以为负；差别来自点电荷自能未包含在
有限的相互作用能中。

由能量求力时必须说明控制量。若系统的几何参数为 $x$，在电荷保持不变的孤立条件下可用
\begin{equation}
 F_x=-\left(\frac{\partial U}{\partial x}\right)_Q.
\end{equation}
若系统接在恒压电源上，电源还会交换电荷并做功，不能只对电容器场能
$CU^2/2$ 机械求导；电容器问题中会进一步区分恒电荷与恒电压两种边界条件。

\subsection*{电势的叠加、梯度与等势面}

静电力是保守力，可定义电势差
\[
 \varphi_B-\varphi_A=-\int_A^B\vec E\cdot\dd\vec l.
\]
电势是单位正电荷的势能，只由场源和位置决定；势能 $U=q\varphi$ 还与被放入的电荷有关。通常
把无穷远取为零势，但无限长线电荷或无限大平面的势在无穷远发散，此时只能选有限参考点讨论
电势差。

点电荷电势为 $q/(4\pi\varepsilon_0r)$，多个点电荷直接作代数和。连续分布为
\[
 \varphi(\vec r)=\frac1{4\pi\varepsilon_0}
 \int\frac{\dd q'}{|\vec r-\vec r'|}.
\]
先求标量电势再取梯度，常比直接积分矢量电场容易：
\[
 \vec E=-\nabla\varphi.
\]
负号表示电场指向电势下降最快的方向。等势面与电场线正交，沿等势面移动电荷时电场不做功；
等势面越密，法向电势变化越快，场强越大。

\subsection*{电荷系统能量与微观能标}

把点电荷逐个从无穷远缓慢搬到指定位置，外力做功储存在电荷系统中：
\[
 U=\frac12\sum_{i\ne j}
 \frac1{4\pi\varepsilon_0}\frac{q_iq_j}{r_{ij}}
 =\frac12\sum_iq_i\varphi_{\rm others}(\vec r_i).
\]
因子 $1/2$ 用来避免每一对相互作用重复计算。连续分布相应写成
$U=\tfrac12\int\rho\varphi\,\dd V$；点电荷自身的发散自能不包含在普通的粒子间相互作用能中。
求系统所受力时，可在保持其他独立坐标不变的条件下使用 $F_x=-\partial U/\partial x$。

电子伏特是能量单位：带一个元电荷的粒子通过 $1\,\mathrm V$ 电势差，能量改变
$1\,\mathrm{eV}=1.602\times10^{-19}\,\mathrm J$。原子、核与高能过程跨越的尺度很大，用 eV、keV、
MeV 表示比焦耳直观。电势高低本身不能直接判断带负电粒子的运动方向，应始终从
$\Delta K=-\Delta U=-q\Delta\varphi$ 判断。

氢原子的经典圆轨道模型把库仑力作为向心力，再配合 Bohr 角动量量子化，可得到能级按
$-1/n^2$ 分布；Rutherford 散射则用正电原子核对 $\alpha$ 粒子的库仑排斥解释大角偏转。经典轨道
图像不能替代量子理论，但它清楚展示了电势能、动能和特征尺度如何由基本常量组合出来。

在远离局域电荷分布处，电势也可作多极展开。总电荷不为零时最低阶为 $Q/r$；总电荷为零但偶极矩
不为零时，最低阶为 $\vec p\cdot\hat r/r^2$；若偶极矩也为零，四极项才主导。每高一阶都多一个
$1/r$ 的衰减，因此观察距离越远，越只看得到最低非零多极矩。


\subsection*{从电场力做功到电势}

静电力是保守力。把试探电荷 $q$ 从 $A$ 移到 $B$，电场力做功只由始末位置决定，定义
\[
  W_{A\to B}=q(\varphi_A-\varphi_B),
  \qquad U=q\varphi.
\]
电势是单位正电荷的势能，是描述电场的标量函数；静电势能则属于具体电荷与外电场组成的系统。
同一点的电势与放入多大试探电荷无关，势能却正比于试探电荷。若试探电荷为负，它受力方向与
电场相反，沿电场方向移动时势能反而增加。因此判断能量变化时不能只背“沿电场电势降低”，还要
同时考虑电荷符号。

电势零点可以任意选择，真正可测的是电势差。孤立有限电荷分布常取无穷远电势为零；无限长线
电荷或无限大带电平面的绝对电势会发散，只能选择有限参考点讨论电势差。多个点电荷产生的电势
直接作代数叠加：
\[
  \varphi(\vect r)=\frac1{4\pi\varepsilon_0}
  \sum_i\frac{q_i}{|\vect r-\vect r_i|}.
\]
因为电势是标量，复杂电荷分布往往先求电势再求梯度，比直接叠加电场矢量更简洁。

电场与电势的局部关系为
\[
  \vect E=-\nabla\varphi.
\]
负号表示电场指向电势下降最快的方向。直角坐标中三个分量分别是
$E_x=-\partial\varphi/\partial x$、$E_y=-\partial\varphi/\partial y$、
$E_z=-\partial\varphi/\partial z$。等势面上任意切向位移满足 $\dd\varphi=0$，所以电场只能
沿法向。图30.1把匀强电场、点电荷和偶极场并列，是为了同时展示平行面、同心球面和复杂曲面
三种几何，而不是只记住“等势线互相平行”这一种特殊情况。

\subsection*{典型分布的电势与等势面几何}

球对称电荷分布在球外的电势等效于总电荷集中在球心。均匀带电球壳内部电场为零，因此内部电势
处处等于球面电势，并不等于零。均匀带电实心球内部电场随 $r$ 线性变化，积分后电势是 $r$ 的
二次函数。计算时应先保证电势在分界面连续，再用无穷远零点确定常量。

两个异号点电荷的零电势面满足到两电荷距离之比为常量，是 Apollonius 球面。当 $+Q$ 位于
$x=a$、$-q$ 位于原点且 $Q>q$ 时，球心在原点左侧，球面包围较小的负电荷；不可能把整个球画在
两电荷之间。若 $Q=q$，距离比变为1，球面退化为中垂平面。图形必须与配方所得球心和半径一致。

两根异号无限长线电荷的横截面问题同样给出 Apollonius 圆。等势圆应包围对应的线电荷，电势
越接近线电荷，圆越小且越靠近该电荷。电场线则是通过两根导线的一族圆弧，并与等势圆正交。
如果画出的等势圆既没有包围正线电荷也没有包围负线电荷，即使形式上有许多圆，位置关系仍是
错误的。

有限长带电线段的电势可由
\[
  \varphi=k\lambda\int_{-a}^{a}
  \frac{\dd x'}{\sqrt{(x-x')^2+\rho^2}}
\]
求得。结果只依赖场点到两端点的距离组合，等势面为绕线段轴旋转的椭球面，配套的正交曲面为
双曲面。这种焦点几何与引力轨道中的焦点概念不同，但都要求图中的端点、焦点和距离定义严格
一致。

\subsection*{电子伏特、原子电离与气体击穿}

电子伏特把电势差和微观粒子能量直接连接起来：一个元电荷通过1伏电势差，能量改变1电子伏特。
原子的第一电离能通常为数电子伏特到二十几电子伏特，并随原子序数呈明显周期性。稀有气体电子
壳层较稳定，第一电离能形成峰值；碱金属最外层电子较易移走，形成低谷。图30.3(a)保留数值纵轴和
元素序列，才能让读者看到“多数处于 $5$--$25\,\mathrm{eV}$”并非随意估计。

核过程的典型能量进入 $\mathrm{MeV}$ 范围，高能加速器则使用 $\mathrm{GeV}$、$\mathrm{TeV}$。
加速器中的电场负责改变粒子能量，磁场主要负责偏转和聚焦。粒子在环形轨道中反复通过加速腔，
每次获得 $q\Delta\varphi$，许多次累积后达到高能；不能把环形隧道本身误解为始终沿切向做功的
静电场。

气体中电子在两次碰撞之间的平均自由程为 $l$，电场在这段距离内提供的能量约为 $eEl$。当它
超过原子激发能时可能发光，超过电离能时会产生新的电子和离子，电子数在连续碰撞中倍增，最终
形成击穿。闪电的明亮分支是电离后的导电通道，不是预先存在的电场线。压强、气体成分、自由程
和电极曲率都会改变击穿条件，所以 $eEl\sim10\,\mathrm{eV}$ 只能作为数量级判断。

\subsection*{电荷系统能量与电场能量}

把多个点电荷从无穷远缓慢搬运到给定位置，外力做功等于系统静电势能：
\[
  U=\sum_{i<j}\frac{q_iq_j}{4\pi\varepsilon_0r_{ij}}
   =\frac12\sum_i q_i\varphi_i^{(\mathrm{others})}.
\]
因子 $1/2$ 用来避免每一对相互作用被计算两次。连续分布形式
$U=\tfrac12\int\rho\varphi\,\dd V$ 也有同样含义。若把自能发散的理想点电荷包含在积分中，
还需明确是否剔除每个电荷与自身的相互作用。

电场本身具有能量，真空中的能量密度为
\[
  u_E=\frac12\varepsilon_0E^2,
  \qquad U=\int u_E\,\dd V.
\]
对均匀带电球壳，既可以用组装电荷的方法计算，也可以把球外电场能量积分；两种结果相同。后者
显示能量储存在有电场的空间，而不是只“附着”在电荷上。由能量对几何参数求导可以得到静电力，
但必须先说明控制条件：孤立系统通常保持电荷不变，接电源系统可能保持电势不变，两种条件下的
能量函数和外力功并不相同。

\subsection*{从电场积分到电势分布}

选定零势面以后，任意点电势由
\[
 \varphi(B)-\varphi(A)=-\int_A^B\vect E\cdot\dd\vect l
\]
确定。静电场中积分与路径无关，所以可以选择计算最方便的路线；若沿等势面移动，线元始终与
电场垂直，积分为零。反过来，等势面上的电势不变意味着电场没有切向分量，故电场线必与等势面
正交。画图时，场线与等势线若以斜角相交，便不是风格问题，而是违反
$\vect E=-\nabla\varphi$ 的物理错误。

对点电荷系，取无穷远为零可写成
\[
 \varphi(\vect r)=\frac1{4\pi\varepsilon_0}\sum_i
 \frac{q_i}{|\vect r-\vect r_i|}.
\]
电势是标量，叠加时只作代数加法；电场是矢量，需要逐分量相加。总电势为零的点不一定电场为零，
电场为零的点也不一定电势为零。前者表示各电荷的标量贡献抵消，后者表示电势在该点梯度为零，
两者条件完全不同。

无限长线电荷的绝对电势不能把无穷远作为零点，因为积分呈对数发散；只能规定某个有限参考半径
$r_0$ 的电势，写成
\[
 \varphi(r)-\varphi(r_0)=-\frac{\lambda}{2\pi\varepsilon_0}
 \ln\frac r{r_0}.
\]
两根线电荷组成的等势线由两个距离之比决定，几何上是 Apollonius 圆。每一簇闭合等势圆必须围住
对应的线电荷；若圆心、半径画错，使电荷落在圆外，便与距离比的几何定义矛盾。

\subsection*{势能、外力功与电场力功}

电荷 $q$ 在外电势 $\varphi$ 中的势能为 $U=q\varphi$。它从 $A$ 移到 $B$ 时，电场力做功
\[
 W_{\rm e}=q[\varphi(A)-\varphi(B)]=-\Delta U.
\]
若缓慢移动、始末动能不变，外力做功 $W_{\rm ext}=\Delta U$。若粒子自由加速，则
$\Delta K=-\Delta U$。这三种说法对应不同的功，符号不能混用。对电子 $q=-e$，它沿电场方向
移动时电势降低但势能升高；“正电荷趋向低电势”不能直接照搬给负电荷。

由势能曲线判断平衡时，$F_x=-\dd U/\dd x$。平衡点满足 $U'(x_0)=0$；若
$U''(x_0)>0$ 是稳定平衡，若 $U''(x_0)<0$ 是不稳定平衡。对受约束运动，只能沿允许方向取导数；
把三维空间的极值结论未经检查地套到一条轨道上，可能漏掉横向不稳定性。

\subsection*{微观能量尺度的读图方法}

第一电离能曲线不是要求逐个死记数值，而是把电子壳层结构转化为可比较的能量尺度。同一周期中，
核电荷增加总体使电子束缚增强，但亚层填充和电子配对会造成局部起伏；到稀有气体形成峰值，进入
下一周期的碱金属后骤降。图中的横轴必须按元素顺序连续排列，纵轴必须保留电子伏特量级，才能
同时读出周期性和数量级。

把电离能与气体放电联系起来时，还需区分电子一次碰撞获得的能量与外加电压。电压很高并不保证
单次自由程内获得足够能量；压强升高会缩短自由程，电场不均匀又使尖端附近更容易先击穿。因而
宏观电压、几何尺度和微观电离能必须通过 $eEl$ 联系，不能只比较两个表面数值。

电势零点的选择不影响电场，也不影响同一电荷在两点间的势能差；但讨论“某点电势为正还是为负”
时必须先说明参考零点。对有限局域电荷系通常取无穷远为零，对无限线电荷或无限平面则应选择有限
参考位置，只比较电势差。

\subsection*{连续带电体的电势积分}

连续电荷分布的电势由点电荷结果叠加得到：
\[
 \varphi(\vect r)=\frac1{4\pi\varepsilon_0}
 \int\frac{\rho(\vect r')}{|\vect r-\vect r'|}\,\dd V',
\]
线分布和面分布则把 $\rho\dd V'$ 换成 $\lambda\dd l'$、$\sigma\dd S'$。因为电势是标量，
积分通常比直接分解电场容易；求出 $\varphi$ 后再用 $\vect E=-\nabla\varphi$ 得到电场。若场点
落在连续分布内部，要注意积分核在邻近电荷元处是否可积；对理想点电荷，自能部分会发散，不能
未经说明地混入连续分布能量。

均匀带电薄球壳是典型例子。球外电势等于把全部电荷集中在球心：
\[
 \varphi_{\rm out}(r)=\frac1{4\pi\varepsilon_0}\frac Qr.
\]
球内电场为零，所以球内电势处处等于表面值
\[
 \varphi_{\rm in}(r)=\frac1{4\pi\varepsilon_0}\frac QR.
\]
电势在球面连续，法向电场却因面电荷发生跃变。对于均匀带电实心球，利用 Gauss 定律先求球内
场，再从表面积分到内部，可得
\[
 \varphi(r)=\frac1{4\pi\varepsilon_0}\frac{Q}{2R}
 \left(3-\frac{r^2}{R^2}\right),\qquad r\le R.
\]
在 $r=R$ 处，它与球外电势及其一阶导数连续；在球心电势取得最大值但电场为零。

\subsection*{Poisson 方程与电势图像}

把 $\vect E=-\nabla\varphi$ 代入 Gauss 定律的微分形式，得到
\[
 \boxed{\nabla^2\varphi=-\frac{\rho}{\varepsilon_0}}.
\]
无电荷区域满足 Laplace 方程。电势曲面的弯曲程度由当地电荷密度决定，但电势本身还受整个空间
电荷和边界条件影响。一个区域内 $\rho=0$ 并不表示 $\varphi=0$ 或 $\vect E=0$；它只表示电势
在该区域为调和函数。

沿一条电场线，取弧长正方向与 $\vect E$ 同向，有
$\dd\varphi/\dd s=-E<0$，所以电势必沿场线降低。等势线画得越密，表示相同电势差发生在越短
距离内，场强越大；但只有在相邻等势线取相同电势间隔时，线距才能用于比较场强。图示若没有标明
等势值，线的疏密只能作定性提示，不能直接当成定量比例尺。


\Needspace{0.34\textheight}
\begin{origexamplebox}{例题 30.1【电势算电场】}
利用“先求电势，再求电场”的方法，重新计算上一讲中的典型带电系统：点电荷系、有限带电体
和无限长或无限大的电荷分布。分析哪些情况下不能取无穷远为零电势，并给出处理方法。
\end{origexamplebox}

\begin{xesolution}
方法的核心是
\[
  \varphi(\vect r)=k\int\frac{\dd q}{|\vect r-\vect r'|},
  \qquad
  \vect E=-\nabla\varphi,
  \qquad k=\frac1{4\pi\varepsilon_0}.
\]

\textbf{点电荷系。}直接标量叠加
\[
  \varphi=k\sum_i\frac{Q_i}{r_i},
\]
然后求梯度即得
\[
  \vect E=k\sum_i\frac{Q_i}{r_i^3}\vect r_i.
\]
这种方法对偶极子、多极子尤其方便，因为远场展开先在标量电势上进行，再求导即可。

\textbf{均匀带电圆盘。}半径 $R$、面密度 $\sigma$ 的圆盘，在轴线上距圆心 $z>0$ 处的电势为
\[
  \varphi(z)=k\int_0^R\frac{2\pi\sigma s\dd s}{\sqrt{z^2+s^2}}
  =\frac{\sigma}{2\varepsilon_0}
  \bigl(\sqrt{z^2+R^2}-z\bigr).
\]
因此
\[
  \boxed{E_z=-\frac{\dd\varphi}{\dd z}
  =\frac{\sigma}{2\varepsilon_0}
  \left(1-\frac{z}{\sqrt{z^2+R^2}}\right)}.
\]

\textbf{无限长直线。}若线电荷密度为 $\lambda$，$E_r=\lambda/(2\pi\varepsilon_0r)$。
从有限参考半径 $r_0$ 积分得
\[
  \boxed{\varphi(r)-\varphi(r_0)
  =-\frac{\lambda}{2\pi\varepsilon_0}\ln\frac r{r_0}}.
\]
因为 $r\to\infty$ 时对数发散，所以不能令无穷远处电势为零。

\textbf{无限大平面。}若面电荷密度为 $\sigma$，取平面为 $z=0$，则
\[
  E_z=\frac{\sigma}{2\varepsilon_0}\operatorname{sgn}(z),
\]
取 $\varphi(0)=0$ 后
\[
  \boxed{\varphi(z)=-\frac{\sigma}{2\varepsilon_0}|z|}.
\]
它同样不能用无穷远作为零电势。对非局域电荷分布，正确做法是选定有限参考点，始终只计算
电势差；若多个无限分布组合后总场在远处衰减足够快，也可先把组合视为整体再选零点。
\end{xesolution}

\Needspace{0.37\textheight}
\begin{origexamplebox}{例题 30.2【均匀带电球】}
半径为 $R$、总电荷为 $Q$ 的金属薄球壳均匀带电。求：
\begin{enumerate}[label=(\arabic*),leftmargin=2.2em]
  \item 全空间的总电场能量；
  \item 保持球壳不被静电斥力胀开所需的压强；
  \item 若改为半径 $R$、总电荷 $Q$、均匀带电的绝缘实心球，前两问如何变化；
  \item 为什么静电平衡时常见的是均匀带电金属球壳，而均匀体电荷只适用于绝缘体？
\end{enumerate}
\sphereEnergyDiagram
\end{origexamplebox}

\begin{xesolution}
\textbf{(1) 球壳的电场能量。}球内 $E=0$，球外
\[
  E(r)=\frac{Q}{4\pi\varepsilon_0r^2}.
\]
于是
\begin{align*}
  U&=\int_R^\infty\frac12\varepsilon_0E^2\,4\pi r^2\dd r\\
   &=\boxed{\frac{Q^2}{8\pi\varepsilon_0R}}.
\end{align*}
这也等于 $\tfrac12Q\varphi(R)$。

\textbf{(2) 静电压强。}球面外侧的电场为
$E_{\rm out}=Q/(4\pi\varepsilon_0R^2)$，内侧为零。表面电荷受到的平均场应排除电荷元自身的
奇异场，等于 $(E_{\rm out}+E_{\rm in})/2$。故向外静电压强
\[
  \boxed{p=\frac12\varepsilon_0E_{\rm out}^2
  =\frac{Q^2}{32\pi^2\varepsilon_0R^4}}.
\]
外界必须提供同样大小的向内机械压强。

\textbf{(3) 均匀带电绝缘实心球。}体电荷密度
$\rho=3Q/(4\pi R^3)$。高斯定律给出
\[
  E(r)=
  \begin{cases}
    \dfrac{Qr}{4\pi\varepsilon_0R^3},&r<R,\\[5pt]
    \dfrac{Q}{4\pi\varepsilon_0r^2},&r\ge R.
  \end{cases}
\]
把球内与球外电场能量分别积分：
\[
  U_{\rm in}=\frac{Q^2}{40\pi\varepsilon_0R},
  \qquad
  U_{\rm out}=\frac{Q^2}{8\pi\varepsilon_0R},
\]
所以
\[
  \boxed{U_{\rm solid}=\frac{3Q^2}{20\pi\varepsilon_0R}}.
\]
绝缘体内各体积元都受电力，不能用单一“表面压强”完整描述其机械平衡；需要体内随半径变化的
应力分布。其表面处电场仍与球壳相同，但这不足以决定整个实心球的内部应力。

\textbf{(4)} 导体内有自由电荷。静电平衡时导体内部电场必须为零，净电荷只能分布在表面；
球形对称下表面分布均匀。绝缘体中的电荷可以被束缚在体内，才可能维持均匀体电荷分布。
\end{xesolution}

\Needspace{0.3\textheight}
\begin{origexamplebox}{例题 30.3【相互作用能】}
一个系统含有 $10$ 个点电荷，第 $i$ 个电荷为 $Q_i$，任意两电荷间距为 $r_{ij}$。
把这些电荷从无穷远逐个搬到给定位置，外力至少需要做多少功？
\end{origexamplebox}

\begin{xesolution}
准静态组装，且忽略辐射和其他耗散时，外力做功等于系统静电势能的增量。初态各电荷相距
无穷远，取 $U_i=0$，因此
\[
  \boxed{W_{\rm ext}=U_f
  =\frac{1}{4\pi\varepsilon_0}
  \sum_{1\le i<j\le10}\frac{Q_iQ_j}{r_{ij}}}.
\]
也可写成
\[
  \boxed{W_{\rm ext}=\frac12\sum_{i=1}^{10}
  Q_i\varphi_i^{(\mathrm{others})}}.
\]
若结果为正，组装时外力净输入能量；若为负，电场对外界净做功，外力必须吸收能量。答案与
搬运次序无关，因为静电力是保守力。
\end{xesolution}

\Needspace{0.34\textheight}
\begin{origexamplebox}{例题 30.4【零等势面】}
两个点电荷 $+Q$、$-q$ 分别置于 $(a,0,0)$ 和原点，求以无穷远为零电势时的零等势面方程。
\zeroPotentialSphereDiagram
\end{origexamplebox}

\begin{xesolution}
任一点 $P(x,y,z)$ 到 $+Q$ 和 $-q$ 的距离分别为
\[
  r_+=\sqrt{(x-a)^2+y^2+z^2},
  \qquad
  r_-=\sqrt{x^2+y^2+z^2}.
\]
零电势条件为
\[
  k\left(\frac Q{r_+}-\frac q{r_-}\right)=0,
  \qquad
  \frac{r_+}{r_-}=\frac Qq.
\]
若 $Q\ne q$，平方并配方，得到球面
\[
  \boxed{\left(x+\frac{aq^2}{Q^2-q^2}\right)^2+y^2+z^2
  =\left(\frac{aQq}{|Q^2-q^2|}\right)^2}.
\]
当 $Q>q$ 时，它包围原点处的较小负电荷；当 $Q<q$ 时，它包围正电荷。若 $Q=q$，上式的
球半径趋于无穷，零等势面退化为两电荷连线的中垂面 $x=a/2$。
\end{xesolution}

\Needspace{0.36\textheight}
\begin{origexamplebox}{例题 30.5【无限长导线】}
两根与 $z$ 轴平行的无限长直导线，线电荷密度分别为 $+\lambda$ 和 $-\lambda$，它们与
$xOy$ 平面的交点分别为 $(a,0)$ 和 $(-a,0)$。求横截面内任意等势线与电场线的方程。
\twoLineDiagram
\end{origexamplebox}

\begin{xesolution}
设场点到 $+\lambda$、$-\lambda$ 导线的距离分别为
\[
  r_+=\sqrt{(x-a)^2+y^2},
  \qquad
  r_-=\sqrt{(x+a)^2+y^2}.
\]
单根无限长直线的绝对电势不能以无穷远为零，但异号组合的发散常量抵消。可取
\[
  \boxed{\varphi(x,y)=\frac{\lambda}{2\pi\varepsilon_0}
  \ln\frac{r_-}{r_+}}.
\]
令 $K=\exp(2\pi\varepsilon_0\varphi_0/\lambda)$，等势线满足
\[
  \boxed{\frac{r_-}{r_+}=K}.
\]
$K\ne1$ 时展开为圆
\[
  \boxed{\left(x-a\frac{K^2+1}{K^2-1}\right)^2+y^2
  =\left(\frac{2aK}{|K^2-1|}\right)^2};
\]
$K=1$ 时为 $y$ 轴。

电场线与等势线正交。利用二维复势的共轭函数，或直接积分微分方程，可得电场线族
\[
  \boxed{\arg(x-a+\mathrm i y)-\arg(x+a+\mathrm i y)=\psi_0}.
\]
把它写成实方程：令 $C=\cot\psi_0$，则
\[
  \boxed{x^2+y^2-a^2=2aCy}.
\]
这同样是一族圆，圆心位于 $y$ 轴，并且与所有等势圆正交。绕各导线平移到三维后，得到相应
的等势柱面和电场线所在的柱面。
\end{xesolution}

\Needspace{0.36\textheight}
\begin{origexamplebox}{例题 30.6【有限长导线】}
一根均匀带电的有限长直导线沿 $x$ 轴放置，从 $(-a,0,0)$ 延伸到 $(a,0,0)$，线电荷密度为
$\lambda$。系统绕 $x$ 轴旋转对称。求等势面方程及子午面内的电场线方程。
\finiteLineDiagramThirty
\end{origexamplebox}

\begin{xesolution}
用柱坐标 $(x,\rho)$ 描述场点，$\rho=\sqrt{y^2+z^2}$。到两端点 $(a,0)$、$(-a,0)$ 的距离
分别记为
\[
  r_+=\sqrt{(x-a)^2+\rho^2},
  \qquad
  r_-=\sqrt{(x+a)^2+\rho^2}.
\]
电势积分为
\begin{align*}
  \varphi(x,\rho)
  &=k\lambda\int_{-a}^a
  \frac{\dd x'}{\sqrt{(x-x')^2+\rho^2}}\\
  &=\boxed{k\lambda\ln
  \frac{r_++r_-+2a}{r_++r_--2a}}.
\end{align*}
因此等势面满足
\[
  \boxed{r_++r_-=2a\cosh\mu_0},
\]
其中 $\mu_0>0$ 由所选电势决定。这是以导线两端为焦点的旋转椭球面。

引入长旋转椭球坐标
\[
  x=a\cosh\mu\cos\nu,
  \qquad
  \rho=a\sinh\mu\sin\nu.
\]
上式电势只依赖 $\mu$，故电场沿 $\mu$ 方向；与 $\mu=\text{常量}$ 正交的坐标线
$\nu=\nu_0$ 就是子午面中的电场线。因此
\[
  \boxed{r_--r_+=2a\cos\nu_0},
\]
它们是以两端点为焦点的双曲线，绕 $x$ 轴旋转后形成相应的双曲面。电场线从正线电荷出发，
在远处逐渐呈径向分布。
\end{xesolution}

\Needspace{0.36\textheight}
\begin{origexamplebox}{例题 30.7【同性相斥】}
质量为 $m$、电荷量为 $+q$ 的小球 $A$ 由一根不可伸长轻绳悬挂。手持同样带电的小球 $B$
缓慢排斥 $A$，最终 $B$ 位于 $A$ 原来的最低位置，而 $A$ 上升高度 $h$。求此过程中手对
$B$ 做的功以及电场能量的增量。
\repulsionDiagram
\end{origexamplebox}

\begin{xesolution}
设绳长为 $l$、终态绳与竖直方向夹角为 $\theta$。由几何关系
\[
  h=l(1-\cos\theta),
\]
而终态两球距离
\[
  r=2l\sin\frac\theta2.
\]
终态 $A$ 静止，受重力 $mg$、绳张力 $T$ 和沿 $BA$ 方向的库仑斥力 $F_e$。弦 $BA$ 与
水平方向的夹角为 $\theta/2$。沿水平、竖直方向列平衡方程：
\begin{align*}
  T\sin\theta&=F_e\cos\frac\theta2,\\
  T\cos\theta+F_e\sin\frac\theta2&=mg.
\end{align*}
消去 $T$，得到
\[
  \boxed{F_e=2mg\sin\frac\theta2}.
\]
另一方面，几何关系可以写成
\[
  h=2l\sin^2\frac\theta2,
  \qquad
  r=2l\sin\frac\theta2,
\]
所以
\[
  \frac hr=\sin\frac\theta2.
\]
初态两球相距无穷远，静电相互作用能为零；终态相互作用能为
\[
  \Delta U_e=\frac{kq^2}{r}=F_er
  =2mg\sin\frac\theta2\,r.
\]
利用 $r\sin(\theta/2)=h$，立刻得到
\[
  \boxed{\Delta U_e=2mgh}.
\]

过程足够缓慢，始末动能均为零，绳张力不做功。因此手对 $B$ 做功等于系统重力势能与静电
相互作用能的增量：
\[
  \boxed{W_{\rm hand}=mgh+\Delta U_e=3mgh}.
\]
若只问“电场能量的增量”，即为 $2mgh$。电荷量 $q$ 决定实现给定高度 $h$ 时所需的绳长与
偏角，但在几何和平衡条件共同约束下，最终能量结果与 $q$ 消去。
\end{xesolution}

\Needspace{0.34\textheight}
\begin{origexamplebox}{例题 30.8【Bohr 氢原子】}
在经典 Bohr 模型中，电子以光速的 $1/137$ 绕质子作圆周运动。求氢原子的总机械能，并分析
为什么它是负值。
\bohrDiagram
\end{origexamplebox}

\begin{xesolution}
库仑力提供向心力：
\[
  \frac{mv^2}{r}=\frac{ke^2}{r^2}
  \quad\Longrightarrow\quad
  r=\frac{ke^2}{mv^2}.
\]
动能与电势能分别为
\[
  K=\frac12mv^2,
  \qquad
  U=-\frac{ke^2}{r}=-mv^2.
\]
因此总能量
\[
  \boxed{E=K+U=-\frac12mv^2}.
\]
取 $v=c/137$、$m=m_e$，得到
\[
  E=-\frac12m_ec^2\frac1{137^2}
  \approx\boxed{-\SI{13.6}{eV}}.
\]
负号来自能量零点的选择：电子与质子相距无穷远且静止时取 $E=0$。束缚态电子要逃到无穷远，
外界至少要输入 $\SI{13.6}{eV}$，所以束缚态总能量低于零。经典圆轨道不能解释原子的稳定性，
这里的结果属于 Bohr 半经典模型；完整描述需要量子力学。
\end{xesolution}

\Needspace{0.38\textheight}
\begin{origexamplebox}{例题 30.9【$\alpha$ 粒子散射】}
电荷量为 $+2e$、初动能为 $K$ 的 $\alpha$ 粒子射向电荷量为 $+Ze$ 的原子核。把原子核视为
固定，回答：
\begin{enumerate}[label=(\arabic*),leftmargin=2.2em]
  \item 轨迹如何随初始状态变化？
  \item 入射粒子束的微分散射截面是多少？
  \item 历史上谁首先完成了这一实验，实验揭示了什么？
\end{enumerate}
\rutherfordDiagram
\end{origexamplebox}

\begin{xesolution}
记库仑斥力常量
\[
  \kappa=k(2e)(Ze)=2kZe^2,
  \qquad K=\frac12mv_\infty^2.
\]

\textbf{(1) 轨迹。}这是正能量平方反比斥力问题，轨迹为双曲线。正碰时瞄准参数 $b=0$，
粒子沿原路返回，最近距离由能量守恒给出
\[
  \boxed{r_{\min}=\frac{\kappa}{K}}.
\]
一般 $b>0$ 时，角动量 $L=mv_\infty b$ 守恒。散射角 $\theta$ 与瞄准参数的关系为
\[
  \boxed{b=\frac{\kappa}{2K}\cot\frac\theta2}.
\]
瞄准参数越小，散射角越大；$b\to0$ 时 $\theta\to\pi$，$b\to\infty$ 时
$\theta\to0$。双曲线偏心率为
\[
  \epsilon=\sqrt{1+\frac{4K^2b^2}{\kappa^2}}>1.
\]

\textbf{(2) Rutherford 微分截面。}轴对称散射中，入射环带面积
$\dd\sigma=2\pi b|\dd b|$，对应出射立体角
$\dd\Omega=2\pi\sin\theta\,\dd\theta$。因此
\[
  \frac{\dd\sigma}{\dd\Omega}
  =\frac{b}{\sin\theta}\left|\frac{\dd b}{\dd\theta}\right|.
\]
代入 $b=(\kappa/2K)\cot(\theta/2)$，得到
\[
  \boxed{\frac{\dd\sigma}{\dd\Omega}
  =\left(\frac{\kappa}{4K}\right)^2
  \frac1{\sin^4(\theta/2)}}
  =\boxed{\left(\frac{kZe^2}{2K}\right)^2
  \frac1{\sin^4(\theta/2)}}.
\]
截面在小角度处很大，在大角度处迅速减小；少量大角散射事件正是原子核式结构的关键信号。

\textbf{(3) 历史背景。}1909 年，Hans Geiger 和 Ernest Marsden 在 Ernest Rutherford 指导下
完成金箔散射实验；Rutherford 于 1911 年提出核式原子模型并给出散射理论。实验发现少数
$\alpha$ 粒子发生大角度甚至近乎反向散射，表明原子的正电荷和绝大部分质量集中在尺度很小的
原子核中，而原子的绝大部分空间是空的。
\end{xesolution}
